% GENERATED FILE. Edit canonical lessons, then run npm run generate:books.
% canonical-source-hash: fnv1a64:ce32da225d1b73fd
% canonical-lessons: KA-C70-say, KA-C70-quotative, KA-S138-vowel-sign-ee, KA-C70-because, KA-C70-i-think

\chapter{Saying So, and Saying Why}
\label{ch:ka-saying-why}
\hlchaptermodality{70}

\begin{chapteropening}
\textbf{By the end of this chapter:} \emph{I can report what was said, give a reason, and say what I think.}

Complete KA-C70-i-think and give an opinion with the quotative and a reason for it.
\end{chapteropening}

\section[hēḷu]{\kn{ಹೇಳು} (hēḷu) — to say, to tell}

\label{lesson:KA-C70-say}



\begin{quote}
Before the new one: which character did you trace last, and what does \emph{yāke} ask?
\end{quote}

\subsection*{You'll want to know: \kn{ಹೇಳು}}
\textbf{\kn{ಹೇಳು}} (\emph{hēḷu}) — \textquotedblleft{}to say, to tell\textquotedblright{}.

You have had \kn{ಮಾತಾಡು} (\emph{mātāḍu}, \textquotedblleft{}to speak\textquotedblright{}) for a while. The two are not interchangeable: \kn{ಮಾತಾಡು} is the activity of talking, and \kn{ಹೇಳು} is saying a particular thing to a particular person. You speak Kannada; you tell somebody the way.

You have also read the polite command form of it already, in the sentence a stranger needs most: \textbf{\kn{ದಯವಿಟ್ಟು} \kn{ದಾರಿ} \kn{ಹೇಳಿ}} (\emph{dayaviṭṭu dāri hēḷi}) — \textquotedblleft{}please tell me the way.\textquotedblright{} The \textbf{-\kn{ಇ}} on the end is the polite request ending, so \kn{ಹೇಳಿ} is \textquotedblleft{}please tell\textquotedblright{}.

The reason this verb arrives here rather than with the other verbs is what comes next: it is the verb a reported sentence hangs on, and until the following lesson there was nothing to hang.

\begin{grammarlens}[title={What you've built}]
A verb with a hole in it, waiting for something to be said.
\end{grammarlens}

\subsection*{Your turn}
\begin{itemize}
\raggedright
  \item \emph{Say it:} \emph{hēḷu}
  \item \emph{Say it:} \emph{hēḷi} — the polite one
  \item \emph{Say it:} \emph{dayaviṭṭu dāri hēḷi}
  \item \emph{Recall:} say \emph{yāke}, then ask somebody why
  \item \emph{Recall:} say \emph{nanagū gottilla} once more
  \item \emph{Recall:} say \emph{\kn{ೂ}} once more
  \item \emph{Recall:} from much earlier — say \emph{mara}, \emph{mōḍa}, \emph{haḷeya}, \emph{hosa}, and say what each one means
\end{itemize}

\begin{tcolorbox}[breakable,colback=teal!4,colframe=teal!35!black,title={Before you move on}]
What does \kn{ಹೇಳು} mean? (\textquotedblleft{}to say, to tell\textquotedblright{}.) What is the polite form? (\emph{hēḷi}.)
\end{tcolorbox}
\section[anta]{\kn{ಅಂತ} (anta) — that}

\label{lesson:KA-C70-quotative}



\begin{quote}
Before the new one: what does \emph{hēḷu} mean, and what is its polite form?
\end{quote}

\subsection*{You'll want to know: \kn{ಅಂತ}}
\textbf{\kn{ಅಂತ}} (\emph{anta}) — the word that closes off something said or thought.

English opens a report with \emph{that} and lets it run: \textquotedblleft{}he said \textbf{that} there is no coffee.\textquotedblright{} Kannada does the opposite. It says the thing first, exactly as it was said, and then shuts it with \kn{ಅಂತ}:

\textbf{\kn{ಕಾಫಿ} \kn{ಇಲ್ಲ} \kn{ಅಂತ} \kn{ಹೇಳಿ}} (\emph{kāphi illa anta hēḷi}) — \textquotedblleft{}please say that there is no coffee.\textquotedblright{} Literally: \textquotedblleft{}there is no coffee — thus — tell.\textquotedblright{}

Nothing inside the quoted part changes. \kn{ಕಾಫಿ} \kn{ಇಲ್ಲ} is the same three syllables whether you are saying it or reporting it.

\begin{grammarlens}[title={Grammar Lens: the quotation closes before the verb arrives}]
\kn{ಅಂತ} comes after the clause it closes, and the verb of saying or thinking comes after that. So a Kannada report is built back to front from an English one:

\noindent
\begin{tabularx}{\linewidth}{@{}>{\raggedright\arraybackslash}X>{\raggedright\arraybackslash}X>{\raggedright\arraybackslash}X@{}}
\toprule
\textbf{what is said} & \textbf{closer} & \textbf{the verb} \\
\midrule
\kn{ಕಾಫಿ} \kn{ಇಲ್ಲ} & \kn{ಅಂತ} & \kn{ಹೇಳಿ} \\
\bottomrule
\end{tabularx}

That order is why this word had to wait for a verb like \kn{ಹೇಳು}. Until there was something to put in the third column, the first two had nowhere to go.

The written twin of \kn{ಅಂತ} is \textbf{\kn{ಎಂದು}} (\emph{endu}). Same job, older shape, and you will meet it in the very next lesson without being told — it is hiding inside a word you are about to learn.
\end{grammarlens}

\begin{grammarlens}[title={What you've built}]
A whole sentence, used as the object of another one. That is the first piece of Kannada in this book that is bigger than a phrase.
\end{grammarlens}

\subsection*{Your turn}
\begin{itemize}
\raggedright
  \item \emph{Say it:} \emph{anta}
  \item \emph{Say it:} \emph{kāphi illa anta hēḷi}
  \item \emph{Say it:} take any true sentence you know, add \emph{anta}, then add \emph{hēḷi}
  \item \emph{Recall:} say \emph{hēḷu}, then \emph{hēḷi}
  \item \emph{Recall:} say \emph{gottā} once more
  \item \emph{Recall:} from much earlier — say \emph{kombe}, \emph{cikka}, \emph{alli}, \emph{illi}, and say what each one means
  \item \emph{Recall:} read \textbf{\kn{ಅಲ್ಲ}}, then say \emph{hālū illa, nīrū illa}
\end{itemize}

\begin{tcolorbox}[breakable,colback=teal!4,colframe=teal!35!black,title={Before you move on}]
Where does \kn{ಅಂತ} go — before or after the thing being reported? (After.) What is its written twin? (\emph{endu}.)
\end{tcolorbox}
\section[ē]{\kn{ೇ} — one character, met inside words you already say}

\label{lesson:KA-S138-vowel-sign-ee}



\begin{quote}
Before the new one: \kn{ಲ} — and which other letter in this script sounds close to it?

One character this time. It is a sign rather than a letter, and it has been inside the word for \textquotedblleft{}how\textquotedblright{} since the third chapter.
\end{quote}

\subsection*{Script you'll notice: \kn{ೇ}}
\textbf{\kn{ೇ}} — \emph{ē}.

It is a \textbf{vowel sign}. Hang it on a consonant and the built-in \emph{a} becomes a long \textbf{ē}. The short one, \kn{◌ೆ}, you have been reading for a long time; this is its long partner, and Kannada keeps the two apart in meaning as well as in length.

You already say these, and every one of them has \kn{ೇ} somewhere inside it:

\begin{itemize}
\raggedright
  \item \textbf{\kn{ಹೇಳು}} \emph{hēḷu} — to say, from the lesson before last
  \item \textbf{\kn{ಹೇಗೆ}} \emph{hēge} — how
  \item \textbf{\kn{ಮೇಕೆ}} \emph{mēke} — a goat
  \item \textbf{\kn{ಬೇರು}} \emph{bēru} — a root
\end{itemize}

\subsection*{Writing: \kn{ೇ}}
Put your pen on \kn{ೇ} and follow its line. Copy the shape you can see — slowly, and larger than it is printed.

\begin{quote}
This book does not yet tell you \textbf{where to start the character or which way to travel}. That is a real thing, taught with real variation from school to school, and it is not written down here until it can be written down with a source. Copying what is in front of you needs no such source, and it is how the shape gets into your hand in the meantime.
\end{quote}

\subsection*{Your turn}
\begin{itemize}
\raggedright
  \item \emph{Look:} at these words, and find \kn{ೇ} in the ones that have it
\end{itemize}

\begin{quote}
\kn{ಹೇಳು}  ·  \kn{ಹೇಗೆ}  ·  \kn{ಮೇಕೆ}
\end{quote}

\begin{itemize}
\raggedright
  \item \emph{Trace it:} \kn{ೇ} three times, saying \emph{ē} as you finish each one
  \item \emph{Look:} back at any page of this chapter and find \kn{ೇ} once more
\end{itemize}

\begin{tcolorbox}[breakable,colback=teal!4,colframe=teal!35!black,title={Before you move on}]
Which character is this — \kn{ೇ}? What sound does it carry? (\emph{\textbf{ē}}.) Name one word you already say that contains it.
\end{tcolorbox}
\section[ēkendare]{\kn{ಏಕೆಂದರೆ} (ēkendare) — because}

\label{lesson:KA-C70-because}



\begin{quote}
Before the new one: which sign did you trace last, and what does \emph{anta} close off?
\end{quote}

\subsection*{You'll want to know: \kn{ಏಕೆಂದರೆ}}
\textbf{\kn{ಏಕೆಂದರೆ}} (\emph{ēkendare}) — \textquotedblleft{}because\textquotedblright{}.

Take it apart and there is nothing in it you have not already met.

\kn{ಏಕೆ} (\emph{ēke}) is the written twin of \kn{ಯಾಕೆ} — \textquotedblleft{}why\textquotedblright{}. \kn{ಎಂದರೆ} (\emph{endare}) is the written quotative \kn{ಎಂದು} with the same \textbf{-\kn{ದರೆ}} ending that made \kn{ಆದರೆ} out of a verb: \textquotedblleft{}if one says\textquotedblright{}. So the whole word means \textbf{\textquotedblleft{}if one says why\textquotedblright{}} — that is, \textquotedblleft{}as to why\textquotedblright{}.

\textbf{\kn{ಕಾಫಿ} \kn{ಇಲ್ಲ}, \kn{ಏಕೆಂದರೆ} \kn{ಹಾಲು} \kn{ಇಲ್ಲ}} (\emph{kāphi illa, ēkendare hālu illa}) — \textquotedblleft{}there is no coffee, because there is no milk.\textquotedblright{}

A reason can now be asked for and given. Neither was possible in this book an hour ago.

\begin{grammarlens}[title={Grammar Lens: the same ending, for the third time}]
\kn{ಏಕೆಂದರೆ} opens the reason, exactly where English puts \emph{because}, and the reason comes second. That is worth saying plainly because Kannada does not always work that way: most of its subordinate machinery comes BEFORE the main clause, and this borrowed-looking word is one of the few that comes after.

Notice the ending you have now seen three times. \kn{ಆದರೆ}, \kn{ಹಾಗಾದರೆ}, \kn{ಏಕೆಂದರೆ} — all three are little conditional sentences that hardened into joining words, and \textbf{-\kn{ದರೆ}} is the piece doing it in each.
\end{grammarlens}

\begin{grammarlens}[title={What you've built}]
Ask why, and answer it. With \kn{ಮತ್ತು}, \kn{ಅಥವಾ}, \kn{ಆದರೆ} and the joining ending already in hand, a Kannada sentence can now hold two ideas and a reason.
\end{grammarlens}

\subsection*{Your turn}
\begin{itemize}
\raggedright
  \item \emph{Say it:} \emph{ēkendare}
  \item \emph{Say it:} \emph{kāphi illa, ēkendare hālu illa}
  \item \emph{Say it:} \emph{yāke?} — then answer yourself with \emph{ēkendare} and a reason
  \item \emph{Recall:} say \emph{ādare}, then say \emph{hāgādare}, and name the ending both share
  \item \emph{Recall:} say \emph{alvā} once more
  \item \emph{Recall:} from much earlier — say \emph{bēru}, \emph{ele}, \emph{ūṭa}, \emph{hṛdaya}, and say what each one means
  \item \emph{Recall:} say \emph{gottilla}, then read \textbf{\kn{ನನಗೂ} \kn{ಗೊತ್ತಿಲ್ಲ}}
\end{itemize}

\begin{tcolorbox}[breakable,colback=teal!4,colframe=teal!35!black,title={Before you move on}]
What does \kn{ಏಕೆಂದರೆ} mean? (\textquotedblleft{}because\textquotedblright{}.) What are its two halves? (\emph{ēke}, \textquotedblleft{}why\textquotedblright{}, and \emph{endare}, \textquotedblleft{}if one says\textquotedblright{}.)
\end{tcolorbox}
\section[adu nija anta yōcisuttēne]{\kn{ಅದು} \kn{ನಿಜ} \kn{ಅಂತ} \kn{ಯೋಚಿಸುತ್ತೇನೆ} (adu nija anta yōcisuttēne) — I think that is true}

\label{lesson:KA-C70-i-think}



\begin{quote}
Before the new one: how do you give a reason, and which ending do \emph{ādare} and \emph{ēkendare} share?
\end{quote}

\subsection*{You'll want to know: \kn{ಅದು} \kn{ನಿಜ} \kn{ಅಂತ} \kn{ಯೋಚಿಸುತ್ತೇನೆ}}
\textbf{\kn{ಅದು} \kn{ನಿಜ} \kn{ಅಂತ} \kn{ಯೋಚಿಸುತ್ತೇನೆ}} (\emph{adu nija anta yōcisuttēne}) — \textquotedblleft{}I think that is true.\textquotedblright{}

\kn{ಯೋಚಿಸು} (\emph{yōcisu}, \textquotedblleft{}to think\textquotedblright{}) has been in this book since the chapter of mind verbs, and until the last lesson it could not do the one thing thinking is for. A thought is a whole sentence, and Kannada could not hand it one without \kn{ಅಂತ}.

Read it in Kannada order: \textbf{\kn{ಅದು} \kn{ನಿಜ}} (the thought) — \textbf{\kn{ಅಂತ}} (closed off) — \textbf{\kn{ಯೋಚಿಸುತ್ತೇನೆ}} (\textquotedblleft{}I think\textquotedblright{}).

\begin{grammarlens}[title={Grammar Lens: one frame, and the slot you change}]
Only the first slot ever changes:

\textbf{\kn{ಕಾಫಿ} \kn{ಒಳ್ಳೆಯದು} \kn{ಅಂತ} \kn{ಯೋಚಿಸುತ್ತೇನೆ}} — I think coffee is good. \textbf{\kn{ಅದು} \kn{ನಿಜ} \kn{ಅಲ್ಲ} \kn{ಅಂತ} \kn{ಯೋಚಿಸುತ್ತೇನೆ}} — I think that is not true.

An opinion delivered flat is a claim; the same one with \kn{ಅಲ್ವಾ} after it is an invitation.
\end{grammarlens}

\begin{grammarlens}[title={What you've built}]
Agree, disagree, give a reason, say what you think. That is a conversation.
\end{grammarlens}

\subsection*{Your turn}
\begin{itemize}
\raggedright
  \item \emph{Say it:} \emph{adu nija anta yōcisuttēne}
  \item \emph{Say it:} the same thing about coffee — \emph{kāphi oḷḷeyadu anta yōcisuttēne}
  \item \emph{Say it:} disagree with something, then say what you think instead
  \item \emph{Recall:} say \emph{ēkendare}, then give a reason with it
  \item \emph{Recall:} the whole chapter — \emph{hēḷu}, \emph{anta}, \emph{\kn{ೇ}}, \emph{ēkendare}, and this one
  \item \emph{Recall:} say \emph{\kn{ಲ}} once more
  \item \emph{Recall:} from much earlier — say \emph{bīja}, \emph{bāyi}, \emph{mūgu}, \emph{kivi}, and say what each one means
\end{itemize}

\begin{tcolorbox}[breakable,colback=teal!4,colframe=teal!35!black,title={Before you move on}]
Which slot in the frame changes? (The first — the thought.) Which two never move? (\emph{anta} and the verb.)
\end{tcolorbox}
